\chapter{Implementazione}\label{ch:implementazione}
Il capitolo collega le scelte del progetto a porzioni selezionate dei sorgenti
correnti. Le tavole delle classi in appendice~\ref{app:classi}, le sequenze
della sezione~\ref{sec:sequenze} e il modello dati della
sezione~\ref{sec:database} restano i riferimenti d'insieme: qui interessa
capire dove sono applicate le regole e fin dove valgono.
Il capitolo è completo e proposto per revisione degli autori, non già
approvato. Capitoli~4--5 e appendice~C restano provvisori.

\begingroup
\emergencystretch=6em
\renewcommand{\lstlistingname}{Frammento}
\lstset{language=Java,basicstyle=\ttfamily\footnotesize,
  columns=fullflexible,keepspaces=true,showstringspaces=false,
  breaklines=true,breakatwhitespace=false,tabsize=4,gobble=4,upquote=true,
  numbers=left,numberstyle=\tiny\color{black!60},numbersep=8pt,
  frame=lines,rulecolor=\color{black!25},captionpos=b,
  aboveskip=9pt,belowskip=9pt,
  literate={à}{{\`a}}1 {è}{{\`e}}1 {é}{{\'e}}1
    {ì}{{\`i}}1 {ò}{{\`o}}1 {ù}{{\`u}}1
    {…}{{\normalfont\textellipsis}}1}

\section{Domain Model}\label{sec:impl-dominio}
\subsection{Invarianti locali e stato mutabile}
\code{Vehicle} è una classe finale con costruttore privato: la factory
\code{create} controlla i dati prima di costruire l'oggetto. ID, modello,
VIN e anno sono campi finali; targa, chilometri, base del prezzo, modalità
e stato possono cambiare. Il frammento~\ref{lst:vehicle-invariante}
mostra una regola che riguarda due attributi, anziché due controlli
indipendenti.

% Fonte: src/main/java/it/unifi/drivehub/domain/vehicle/Vehicle.java:59-62
\noindent\begin{minipage}{\linewidth}
\begin{lstlisting}[firstnumber=59,caption={\code{Vehicle.java}: normalizzazione della targa e coerenza con i chilometri.},
  label=lst:vehicle-invariante]
        String cleanPlate = (licencePlate != null && !licencePlate.trim().isEmpty()) ? licencePlate.trim().toUpperCase() : null;
        if (cleanPlate == null && kilometers != 0) {
            throw new IllegalArgumentException("Se un veicolo non ha targa, obbligatoriamente i chilometri devono essere 0.");
        }
\end{lstlisting}
\end{minipage}


La stessa coerenza è controllata dai mutatori: per rimuovere la targa occorre
prima portare i chilometri a zero. La factory rifiuta chilometri e prezzo
negativi, anno fuori dall'intervallo 1900--anno corrente, modello o stato
nulli. Il VIN viene ripulito e convertito in maiuscolo; il controllo
effettivo ammette da uno a diciassette caratteri dopo il trim, senza imporre
diciassette caratteri esatti o verificarne l'alfabeto. L'unicità appartiene
al database. Il confronto con l'anno del modello richiede
\code{validateProductionYear}, invocato dal servizio di catalogo, ed è
anche presidiato dai trigger SQL. Il solo \code{setState} rifiuta null:
non codifica un automa di transizioni; sono i servizi a verificare gli stati
ammessi nei singoli comandi.

\code{Request} mantiene finali tipo, ID di richiesta, veicolo e cliente,
acconto e totale;
venditore e scadenza sono aggiornabili. Le due interpretazioni economiche
del campo \code{deposit} dipendono da \code{RequestType}: con \code{SALE}
è l'acconto della vendita al cliente; con \code{PURCHASE} è la stima per
l'acquisto dal cliente, senza totale concordato.

% Fonte: src/main/java/it/unifi/drivehub/domain/operation/Request.java:54-65
\noindent\begin{minipage}{\linewidth}
\begin{lstlisting}[firstnumber=54,caption={\code{Request.java}: forma dell'accordo e precisione degli importi.},
  label=lst:request-accordo]
        if (type == RequestType.SALE) {
            if (agreedTotal == null || agreedTotal.signum() < 0
                    || deposit.compareTo(agreedTotal) > 0) {
                throw new IllegalArgumentException("Totale concordato e acconto non validi per l'acquisto.");
            }
            // Gli importi dell'accordo devono essere già espressi in centesimi, come l'addebito.
            agreedTotal = agreedTotal.setScale(2, RoundingMode.UNNECESSARY);
            deposit = deposit.setScale(2, RoundingMode.UNNECESSARY);
        } else if (agreedTotal != null) {
            throw new IllegalArgumentException("Il totale concordato riguarda solo gli acquisti del cliente.");
        }
        return new Request(id, vehicleId, clientId, salesmanId, deposit, agreedTotal, expDate, type);
\end{lstlisting}
\end{minipage}


I controlli precedenti a questo estratto rifiutano acconto/stima negativi
e ID di veicolo o cliente nulli. \code{UNNECESSARY} accetta solo importi
rappresentabili in centesimi senza arrotondamento; non corregge un valore
più preciso. \code{getRemainingAmount} sottrae l'acconto dal totale solo
per \code{SALE}. La ricostruzione tramite \code{create} accetta una
scadenza già passata: serve a leggere richieste persistite da elaborare.
È il mutatore \code{setExpDate} a rifiutare una nuova scadenza passata;
\code{isExpired} confronta la data con oggi, escluso oggi stesso.

\subsection{Gerarchie, identità e riferimenti}
La gerarchia \code{User}/\code{AbstractUser} condivide identificativo,
nome, cognome, email e hash della password; \code{Client}, \code{Salesman}
e \code{Manager} specializzano ruolo e apertura della dashboard.
Nome e cognome sono finali, email e hash aggiornabili con validazione.
\code{Client} contiene un riferimento opzionale e sostituibile a
\code{BankAccount}, caricato anche dopo l'autenticazione: non è una
composizione UML deducibile dal solo campo. La gerarchia degli sconti
riusa nel costruttore di \code{Discount} il controllo inclusivo 0--50;
\code{ModelDiscount} e \code{VehicleDiscount} conservano l'ID del bersaglio
e una percentuale finale.

I legami persistenti rimangono prevalentemente scalari:
\code{Vehicle.modelId}, gli ID in \code{Request}, \code{Rent},
\code{Sale} e \code{TestDrive}, e \code{BankAccount.clientId} non
caricano automaticamente gli oggetti collegati. Un ID nullo può indicare
un'entità ancora da inserire; \code{save} restituisce un nuovo oggetto con
l'ID assegnato. Non esistono un'identity map o il rilevamento automatico
delle modifiche: mutare una copia Java non aggiorna il database. Anche
un record di lettura che contiene un \code{Vehicle} non rende immutabile
quel veicolo. Le relazioni mostrate nelle tavole CL vanno lette con questa
distinzione fra ID e riferimenti effettivi.

\section{Business Layer}\label{sec:impl-business}
\subsection{Calcolo dei prezzi e accordo economico}
I servizi coordinano i casi d'uso e le scritture; il calcolatore non esegue
SQL. Il catalogo e le modifiche commerciali competono a
\code{CatalogService}; le consultazioni delle operazioni sono raccolte
in \code{OperationQueryService}. I servizi economici separano vendita
al cliente, acquisto dal cliente e noleggio. Il calcolatore
\code{PriceCalculatorService} centralizza le formule.
Seguendo l'ADR~0001, la modalità manuale seleziona una base alternativa,
non un totale finale imposto.

% Fonte: src/main/java/it/unifi/drivehub/domain/vehicle/Vehicle.java:122-124
\noindent\begin{minipage}{\linewidth}
\begin{lstlisting}[firstnumber=122,caption={\code{Vehicle.java}: selezione della base, prima della formula.},
  label=lst:base]
    public BigDecimal getSelectedBasePrice(Model model) {
        return priceMode == PriceMode.MANUAL ? price : model.getReleasePrice();
    }
\end{lstlisting}
\end{minipage}


Il consigliato parte sempre da \code{Model.releasePrice}; l'effettivo
passa la base selezionata alla stessa formula. Lo sconto promozionale è
il massimo fra modello e veicolo. Seguono, in cascata, fattore 0,80 per
l'usato con targa, $0{,}95^{anni}$ e $0{,}95^{km/20000}$, con anni non
negativi rispetto all'anno corrente. La base manuale è arrotondata ai
centesimi con \code{HALF_UP}. Il risultato della formula è anch'esso
arrotondato ai centesimi, ma fattori e potenze usano \code{double}:
\code{BigDecimal} non rende esatto ogni passaggio del calcolo.
L'acconto è il 20\% del prezzo; la tariffa iniziale giornaliera è lo
0,1\%. Il costo del noleggio usa la somma geometrica con fattore 0,95,
conteggio inclusivo dei giorni e arrotondamento finale: non è la tariffa
giornaliera già arrotondata moltiplicata per i giorni.

% Fonte: src/main/java/it/unifi/drivehub/service/PurchaseService.java:74-82
\noindent\begin{minipage}{\linewidth}
\begin{lstlisting}[firstnumber=74,caption={\code{PurchaseService.java}: prezzo corrente e acconto prima della persistenza dell'accordo.},
  label=lst:letture-prezzo]
            Vehicle vehicle = vehicleDao.findById(vehicleId, conn)
                    .orElseThrow(() -> new IllegalStateException("Veicolo non trovato: " + vehicleId));
            Model model = modelDao.findById(vehicle.getModelId())
                    .orElseThrow(() -> new IllegalStateException("Modello non trovato: " + vehicle.getModelId()));

            Integer mDisc = discountDao.findModelDiscount(model.getId()).map(ModelDiscount::getValue).orElse(0);
            Integer vDisc = discountDao.findVehicleDiscount(vehicle.getId()).map(VehicleDiscount::getValue).orElse(0);
            BigDecimal currentPrice = priceCalculator.calculateVehiclePrice(model, vehicle, mDisc, vDisc);
            BigDecimal deposit = priceCalculator.calculatePurchaseDeposit(currentPrice);
\end{lstlisting}
\end{minipage}


Prima di questo estratto, un UPDATE condizionale porta il veicolo da
\code{SALE} a \code{BLOCKED}; se nessuna riga cambia, l'avvio fallisce.
Subito dopo, il servizio blocca il conto, addebita l'acconto e salva
\code{currentPrice} come \code{agreed_total}. Il saldo rilegge la richiesta
sotto lock, verifica approvazione e scadenza e usa totale meno acconto,
senza ricalcolare il listino. L'esempio 20.000/4.000/16.000 euro resta valido
anche dopo una modifica dei prezzi. Il parametro \code{conn} compare
nella lettura del veicolo ma manca in quelle di modello e sconti:
queste aprono connessioni autonome. La transazione protegge le scritture
dell'accordo, senza offrire uno snapshot comune dell'intero listino.

\subsection{Transazioni, lock e operazioni economiche}
\code{DatabaseManager} apre una connessione per ogni esecuzione
transazionale. Il frammento~\ref{lst:transazione} mostra il percorso con
risultato; quello senza risultato applica lo stesso schema.

% Fonte: src/main/java/it/unifi/drivehub/dao/DatabaseManager.java:88-106
\noindent\begin{minipage}{\linewidth}
\begin{lstlisting}[firstnumber=88,caption={\code{DatabaseManager.java}: commit, tentativo di rollback e conservazione delle cause.},
  label=lst:transazione]
    public <T> T executeInTransactionWithResult(TransactionFunction<T> function) {
        try (Connection conn = getConnection()) {
            conn.setAutoCommit(false);
            try {
                T result = function.execute(conn);
                conn.commit();
                return result;
            } catch (Exception e) {
                try {
                    conn.rollback();
                } catch (SQLException rollbackEx) {
                    e.addSuppressed(rollbackEx);
                }
                throw new RuntimeException("Transazione fallita: eseguito rollback.", e);
            }
        } catch (SQLException e) {
            throw new RuntimeException("Errore di connessione al database durante la transazione.", e);
        }
    }
\end{lstlisting}
\end{minipage}


Il valore torna al chiamante dopo \code{commit}; il try-with-resources
chiude la connessione. Le eccezioni del corpo o del commit raggiungono
il ramo di rollback, poi sono incapsulate con la causa. Un errore del
rollback è aggiunto come eccezione soppressa: il testo del messaggio
non dimostra che il rollback sia riuscito. Il codice intercetta
\code{Exception}, non ogni \code{Error}; un guasto di rete durante il
commit non permette di dedurne sempre l'esito dal client. Non vi sono
retry automatici o transazioni annidate implicite.

La connessione ricevuta deve essere passata realmente ai DAO e alle
callback. Avere lo stesso \code{DatabaseManager} non basta: ogni sua
\code{getConnection} apre una connessione JDBC distinta. Nel noleggio,
anche le letture di modello e sconti ricevono \code{conn}; ciò le fa
partecipare alla transazione, senza impostare un isolamento serializzabile
o congelare ogni lettura concorrente. Il servizio blocca veicolo,
cliente e conto, rilegge la disponibilità, ricalcola il costo e addebita
l'intero importo. I vincoli di esclusione sulle date completano questi
protocolli specifici; non sono una prova di assenza generale di deadlock.

Il registro \code{Sale} conserva importi non negativi e tipo esplicito:
\code{SALE} per la vendita al cliente, \code{PURCHASE} per l'acquisto dal
cliente, \code{PENALTY} per la penale. È distinto dalla richiesta ancora
aperta. Al saldo si registra il totale, si elimina la richiesta e il
veicolo passa a \code{SOLD}. Il rifiuto della vendita al cliente rimborsa
l'acconto; la scadenza trattiene il 20\% dell'acconto, non del totale,
e restituisce il resto. I candidati scaduti sono cercati all'apertura
della dashboard e riletti uno per transazione, senza un job periodico.

Nell'acquisto dal cliente, la proposta salva veicolo \code{WAITING} e
stima, senza pagamento. L'accettazione blocca richiesta, veicolo e conto,
accredita la stima persistita e registra \code{PURCHASE}; il rifiuto
elimina il nuovo veicolo solo se privo di storico, altrimenti ripristina
\code{SOLD}. Il controllo sul VIN noto richiede anche l'ultima vendita al
proponente: avere una qualsiasi operazione associata non prova titolarità.
Il cambio di destinazione ammette solo \code{SALE}/\code{RENT}; per
uscire dal noleggio esclude un noleggio in corso e rimborsa i futuri,
bloccando i conti per ID cliente crescente. Le sequenze
SQ-01--03 documentano le alternative senza doverle ripetere qui.

\subsection{Notifiche sincrone ed errori}
% Fonte: src/main/java/it/unifi/drivehub/event/VehicleEventPublisher.java:28-32
\noindent\begin{minipage}{\linewidth}
\begin{lstlisting}[firstnumber=28,caption={\code{VehicleEventPublisher.java}: gli osservatori ricevono la connessione del chiamante.},
  label=lst:observer]
    public void publishCommitment(VehicleCommitmentEvent event, Connection conn) {
        for (VehicleCommitmentObserver observer : observers) {
            observer.onVehicleCommitted(event, conn);
        }
    }
\end{lstlisting}
\end{minipage}


\code{AppContext} registra \code{TestDriveCancellationObserver}.
Vendita al cliente completata e noleggio confermato pubblicano l'evento
prima del commit. L'osservatore elimina i test drive con data e ora
successive a \code{CURRENT_TIMESTAMP} per il veicolo dell'evento, anche
fuori dal periodo noleggiato, usando la
stessa connessione. Non c'è un thread degli eventi: una sua eccezione
risale dal ciclo e impedisce di proseguire verso il commit, annullando
anche addebito e registrazioni nel normale percorso di rollback.
La figura~\ref{fig:sqob} rende visibile questo confine.
Le chiamate successive agli altri osservatori non sono garantite dopo
un errore; non ci sono coda persistente, outbox o consegna distribuita.

I DAO traducono \code{SQLException} in eccezioni con causa; i servizi
sollevano errori per stati incompatibili e dati mancanti. Alcuni, come
\code{RentalService}, estraggono dal wrapper transazionale la causa di
dominio per mostrarla alla vista; questa politica non è uniforme.
Un limite concreto è l'importo zero: richieste e registro lo ammettono,
ma \code{BankAccount.deposit}/\code{withdraw} richiedono un importo
strettamente positivo. I servizi non saltano questi movimenti: una stima
zero può essere proposta ma fallire all'accettazione. La discrepanza è
segnalata separatamente nelle verifiche, senza correzioni al software.

\subsection{Identità e conto dimostrativo}
L'autenticazione dipende dal contratto \code{PasswordHasher}.
Il contesto sceglie PBKDF2-HMAC-SHA256; ogni hash usa un sale casuale
e incorpora i parametri. Il confronto finale usa \code{MessageDigest.isEqual}.
Viene persistito l'hash, senza conservare la password in chiaro.
Il login legge l'identità dalla vista SQL unificata, verifica l'hash e,
per il cliente, carica il conto separatamente.

% Fonte: src/main/java/it/unifi/drivehub/service/AuthService.java:57-64
\noindent\begin{minipage}{\linewidth}
\begin{lstlisting}[firstnumber=57,caption={\code{AuthService.java}: registrazione del cliente e conto iniziale nella stessa transazione.},
  label=lst:registrazione]
        return db.executeInTransactionWithResult(conn -> {
            Client savedClient = userDao.saveClient(clientToSave, conn);
            String generatedIban = generateUniqueIban();
            BankAccount initialAccount = BankAccount.createInitial(generatedIban, savedClient.getId());
            bankAccountDao.save(initialAccount, conn);
            savedClient.setBankAccount(initialAccount);
            return savedClient;
        });
\end{lstlisting}
\end{minipage}


Il nome e l'hash sono validati prima del corpo mostrato. Il conto nasce
con saldo zero; la ricarica positiva blocca e aggiorna la riga del conto.
Il formato dell'IBAN è controllato sintatticamente: generazione casuale
e vincolo UNIQUE non costituiscono un conto bancario reale né un controllo
bancario della cifra di verifica. Email e password richiedono conferma
della password corrente, ma la verifica precede la transazione di scrittura.
La cancellazione blocca l'identità con \texttt{FOR NO KEY UPDATE}, blocca
il conto del cliente e verifica richieste/storico prima di eliminare;
le FK preservano le attribuzioni e il conto viene eliminato per CASCADE,
anche con saldo positivo. Non è un trasferimento bancario.

Dashboard separate e parametri ID/ruolo non equivalgono a un controllo
generale dell'autorizzazione nei servizi. Inoltre la sequenza utenti
condivisa assegna i nuovi ID, ma le PK restano locali alle tre tabelle;
la ricerca per email ignora le maiuscole mentre UNIQUE e trigger SQL
confrontano i valori esatti. Non si attribuiscono unicità globale o
esclusione di ogni gara concorrente a questi soli controlli.

\section{Presentation Layer}\label{sec:impl-presentazione}
\subsection{Navigazione e responsabilità delle viste}
\code{App} mantiene lo Stage e sostituisce la Scene per login e dashboard.
Passa alle viste servizi, utente e callback di logout. Dopo il login,
invoca \code{user.openDashboard} con un \code{UserDashboardNavigator}
anonimo che costruisce le tre dashboard: la scelta è distribuita fra
tipo dell'utente e overload del navigator.

% Fonte: src/main/java/it/unifi/drivehub/domain/user/Client.java:32-35
\noindent\begin{minipage}{\linewidth}
\begin{lstlisting}[firstnumber=32,caption={\code{Client.java}: selezione dell'overload del navigator con il tipo concreto.},
  label=lst:navigazione]
    @Override
    public void openDashboard(UserDashboardNavigator navigator) {
        navigator.showClientDashboard(this);
    }
\end{lstlisting}
\end{minipage}


Le viste costruiscono controlli e dialoghi, raccolgono input e invocano i
servizi. I diagrammi di navigazione della sezione~\ref{sec:navigazione}
restano validi; questo frammento spiega il meccanismo che li realizza.
Il logout richiama la costruzione del login, senza una sessione utente
globale conservata nel contesto. Alcuni costruttori di comodo delle viste
accedono però direttamente al Singleton: l'iniezione non è uniforme.

\subsection{Caricamenti, righe preparate e ciclo di vita}
Le tabelle del catalogo ricevono \code{CatalogService.VehicleRow}: modello,
base, consigliato, effettivo e tariffa sono preparati dal servizio usando
letture collettive e mappe per ID. Le celle di queste tabelle leggono
le righe già preparate, senza query durante il rendering. Non è una cache
persistente né uno snapshot transazionale comune; altre consultazioni,
come le proprie operazioni e alcune descrizioni di veicolo, usano ancora
letture sincrone.

\code{AsyncLoad} associa al controllo una sorgente \code{Supplier} e una
callback \code{Consumer}. Ogni \code{load} cancella il Task precedente,
cattura la sorgente corrente, crea un nuovo Task e disabilita il controllo
durante il caricamento. \code{call} esegue il Supplier sul thread virtuale
avviato in fondo al frammento; i gestori di completamento di Task sono
eseguiti sul thread JavaFX.

% Fonte: src/main/java/it/unifi/drivehub/presentation/view/AsyncLoad.java:60-73
\noindent\begin{minipage}{\linewidth}
\begin{lstlisting}[firstnumber=60,caption={\code{AsyncLoad.java}: pubblicazione del risultato corrente e stato di errore.},
  label=lst:async]
        task.setOnSucceeded(event -> {
            if (current != task) return;
            apply.accept(task.getValue());
            progress.setVisible(false);
            status.setText("Aggiornato");
            target.setDisable(false);
        });
        task.setOnFailed(event -> {
            if (current != task) return;
            progress.setVisible(false);
            status.setText("Caricamento non riuscito. Riprova.");
            refresh.setText("Riprova");
        });
        Thread.ofVirtual().name("drivehub-table-load").start(task);
\end{lstlisting}
\end{minipage}


L'identità \code{current != task} scarta una risposta superata dopo un
nuovo caricamento. Il successo sostituisce le righe e riabilita il
controllo; il fallimento mostra un messaggio generico e lascia disponibile
\emph{Riprova}, mantenendo il controllo destinazione disabilitato.
Il listener sulla proprietà \code{scene} cancella il Task e azzera
\code{current} quando il controllo viene staccato da una Scene.
Questo presidia la sostituzione della vista, senza costituire un gestore
generale del ciclo di vita di ogni dialogo. La cancellazione del Task
non garantisce l'interruzione immediata di una chiamata JDBC già avviata.

\subsection{Aggiornamenti ed esiti dei comandi}
Per esempio, \code{handlePurchaseAction} mostra un preventivo e chiede
conferma; il servizio ricalcola poi l'accordo. Al ritorno aggiorna saldo,
catalogo e richieste e mostra totale e acconto effettivamente persistiti.
Il saldo rilegge ancora la richiesta nel servizio, anche se i controlli
preventivi della vista usano una copia precedente. Le scritture non
dipendono quindi dall'attualità della riga selezionata nella tabella.

Questi handler restano sincroni sul thread grafico: \code{AsyncLoad}
non rende asincrono ogni gesto dell'applicazione. Anche
\code{Platform.runLater}, usato per il controllo iniziale delle scadenze,
rinvia lavoro sul thread JavaFX, senza spostarlo su un worker. Quel
controllo stampa l'errore su stderr; non applica lo stato Riprova delle
tabelle. I dialoghi intercettano gli errori dei comandi e li presentano
con \code{UiUtils}. Il commit precede il successo mostrato, ma gli
aggiornamenti UI non fanno parte della transazione: una lettura di
aggiornamento fallita dopo il commit può produrre un messaggio di errore
pur avendo già completato l'operazione economica.

\section[Persistenza: livello DAO e mapping oggetti/tabelle]{\raggedright Persistenza: livello DAO e mapping oggetti/tabelle}\label{sec:impl-dao}
\subsection{Contratti JDBC e proprietà della connessione}
Le interfacce DAO espongono operazioni sui tipi del dominio; le
implementazioni \code{Postgres*Dao} contengono SQL e mapping espliciti.
Non c'è un ORM: i contratti stessi conoscono \code{java.sql.Connection}.
Il criterio per il riuso è verificabile nei corpi, non nel nome del DAO.

% Fonte: src/main/java/it/unifi/drivehub/dao/postgres/PostgresVehicleDao.java:61-72
\noindent\begin{minipage}{\linewidth}
\begin{lstlisting}[firstnumber=61,caption={\code{PostgresVehicleDao.java}: lettura autonoma e overload con connessione ricevuta.},
  label=lst:dao-connection]
    public Optional<Vehicle> findById(Long id) {
        try (Connection conn = db.getConnection()) {
            return findById(id, conn);
        } catch (SQLException e) {
            throw new RuntimeException("Errore di connessione durante la ricerca del veicolo id: " + id, e);
        }
    }

    @Override
    public Optional<Vehicle> findById(Long id, Connection conn) {
        return findById(id, conn, false);
    }
\end{lstlisting}
\end{minipage}


Il primo overload apre e chiude la connessione; il secondo passa quella
ricevuta al metodo comune, che chiude solo statement e result set.
La responsabilità di commit e chiusura resta al chiamante. Le varianti
\code{findByIdForUpdate} aggiungono \texttt{FOR UPDATE} e rifiutano
autocommit attivo. Una lettura autonoma chiamata dentro una lambda
transazionale non si unisce per questo alla connessione della lambda.

% Fonte: src/main/java/it/unifi/drivehub/dao/postgres/PostgresVehicleDao.java:195-205
\noindent\begin{minipage}{\linewidth}
\begin{lstlisting}[firstnumber=195,caption={\code{PostgresVehicleDao.java}: transizione condizionata e conteggio delle righe.},
  label=lst:update-condizionale]
    public boolean updateStateAtomically(Long vehicleId, VehicleState expectedOldState, VehicleState newState, Connection conn) {
        String sql = "UPDATE vehicle SET state = ?::vehicle_state WHERE id = ? AND state = ?::vehicle_state";
        try (PreparedStatement ps = conn.prepareStatement(sql)) {
            ps.setString(1, newState.name());
            ps.setLong(2, vehicleId);
            ps.setString(3, expectedOldState.name());
            int affected = ps.executeUpdate();
            return affected > 0;
        } catch (SQLException e) {
            throw new RuntimeException("Errore durante l'aggiornamento atomico dello stato del veicolo.", e);
        }
\end{lstlisting}
\end{minipage}


I placeholder parametrizzano i valori; il servizio decide cosa fare del
booleano restituito. Nel saldo e nell'avvio un risultato falso è un errore,
non un aggiornamento da ignorare. Altri percorsi verificano un conteggio
esattamente pari a uno, ad esempio eliminazione richiesta e aggiornamento
saldo; non ogni UPDATE del progetto applica lo stesso controllo.
L'inserimento usa \texttt{RETURNING id}, verifica che sia restituita una
riga e costruisce l'entità identificata, anziché assegnare un ID al campo
finale dell'oggetto originario.

\subsection{Tipi, nullabilità e ricostruzione}
% Fonte: src/main/java/it/unifi/drivehub/dao/postgres/PostgresRequestDao.java:214-226
\noindent\begin{minipage}{\linewidth}
\begin{lstlisting}[firstnumber=214,caption={\code{PostgresRequestDao.java}: NULL preservati e factory anche in lettura.},
  label=lst:mapping]
    private static Request mapRow(ResultSet rs) throws SQLException {
        Long salesman = rs.getObject("salesman") != null ? rs.getLong("salesman") : null;
        Date expDate = rs.getDate("exp_date");
        return Request.create(
                rs.getLong("id"),
                rs.getLong("vehicle"),
                rs.getLong("client"),
                salesman,
                rs.getBigDecimal("deposit"),
                rs.getBigDecimal("agreed_total"),
                expDate != null ? expDate.toLocalDate() : null,
                RequestType.valueOf(rs.getString("type_str"))
        );
\end{lstlisting}
\end{minipage}


\code{getLong} da solo non distingue un NULL da zero: per il venditore
opzionale il controllo su \code{getObject} conserva null.
\code{getBigDecimal} conserva i valori decimali di \code{deposit} e
\code{agreed_total}; \code{getDate}/\code{toLocalDate} ricostruiscono
una data senza orario, preservandone l'assenza. I campi obbligatori per FK
sono letti come ID, senza caricare le entità collegate.
La scrittura simmetrica usa \code{setNull} con \code{Types.BIGINT} o
\code{Types.DATE}, \code{Date.valueOf} e \code{setBigDecimal}; per il
totale opzionale passa il BigDecimal anche quando nullo. La ricostruzione
attraverso la factory riapplica i controlli del dominio: un dato incompatibile
può fallire durante la lettura, non viene silenziosamente sanato.

I quattro domini chiusi sono enum Java e tipi PostgreSQL distinti, come
stabilito dall'ADR~0002. I cast JDBC mantengono questa corrispondenza:

\begin{center}\small
\begin{tabular}{lll}
Attributo & Tipo Java & Cast SQL in scrittura \\\hline
\code{vehicle.state} & \code{VehicleState} & \code{?::vehicle_state} \\
\code{request.type} & \code{RequestType} & \code{?::request_type} \\
\code{vehicle.price_mode} & \code{PriceMode} & \code{?::price_mode} \\
\code{sale.type} & \code{SaleType} & \code{?::sale_type}
\end{tabular}
\end{center}

I parametri sono impostati da \code{enum.name()}; in lettura i DAO
convertono con \code{valueOf}, in alcuni SELECT dopo \texttt{::text}
e alias. I valori ammessi sono quelli già riportati in
sezione~\ref{sec:database}: il cast non verifica regole su importi,
accordi o associazioni. Aggiungere o rinominare un valore richiede
coerenza fra Java, SQL e dati; non c'è una conversione automatica nascosta.

La gerarchia degli utenti è ricostruita dalla vista \code{user}, UNION ALL
delle tabelle di ruolo, con uno switch che istanzia il sottotipo e lascia
al servizio il caricamento del conto. Le scritture scelgono la tabella
con \code{resolveTable}, un insieme chiuso di ruoli, mentre i valori
restano parametri. Le figure~\ref{fig:rel02} e~\ref{fig:rel03} descrivono
nullabilità e FK complete: il mapping non aggiunge una PK globale degli
utenti o un vincolo di autorizzazione assente nel DDL. Il caso SQLSTATE
23503 nella cancellazione dell'identità è tradotto in errore di dominio;
gli altri errori conservano la causa SQL. Non c'è una tassonomia uniforme
di eccezioni applicative per tutti i DAO.

\section{Idiomi e Design Pattern}\label{sec:impl-pattern}
La classificazione dipende dalle collaborazioni effettive, riscontrabili
nei frammenti e nella figura~\ref{fig:pkg}. Un'etichetta nei Javadoc
non dimostra un pattern; il Mode \#2 non richiede di raggiungerne un
numero prefissato.

\paragraph{Factory statiche e idiomi.}
\code{Vehicle.create} e le factory nominate di \code{Request} e
\code{Sale} nascondono costruttori e riusano la validazione anche nei
mapper. Manca però un Creator con metodo di creazione ridefinito dai
sottotipi: non è Factory Method GoF. Campi finali, enum, record e
try-with-resources sono idiomi Java. Una classe finale può essere mutabile;
un enum di stati non realizza da solo State GoF.

\paragraph{Observer GoF.}
Publisher, contratto, iscrizione/rimozione e notifica degli osservatori
realizzano la collaborazione GoF. Separano la cancellazione dei test drive
dai servizi economici. Il frammento~\ref{lst:observer} mostra il limite:
notifica sincrona, errori propagati e connessione JDBC esplicita.
Un eventuale osservatore con effetti esterni non riceverebbe automaticamente
rollback per quegli effetti. Non è una coda persistente né una notifica UI
dopo il commit.

\paragraph{Navigazione nello stile Visitor.}
\code{openDashboard} seleziona l'overload per il tipo concreto
(frammento~\ref{lst:navigazione}); \code{App} fornisce il navigator.
Questo double dispatch separa dalle entità la costruzione delle dashboard.
Il contratto enumera però i tre ruoli e \code{User} conosce la navigazione,
pur senza importare JavaFX. Un nuovo ruolo modifica navigator, App e mapper:
non basta aggiungere una classe. Non è un Visitor per attraversare
arbitrariamente il grafo del dominio.

\paragraph{DAO e dependency injection.}
DAO è un pattern di accesso ai dati esterno al catalogo GoF: qui isola SQL
e mapping ma espone Connection, senza indipendenza da JDBC o repository
di aggregati DDD. La DI manuale è una tecnica di composizione, anch'essa
esterna ai GoF: \code{AppContext} costruisce e passa i collaboratori.
Riceve il gestore ma sceglie internamente PostgreSQL e PBKDF2; le viste
ricevono servizi concreti. L'inversione delle dipendenze è quindi parziale.

\paragraph{Singleton e composizione.}
Il frammento~\ref{lst:singleton} realizza l'accesso Singleton GoF al contesto
predefinito. Il costruttore pubblico con \code{DatabaseManager} consente
però contesti separati: non impone un'unica istanza nel processo.
La sincronizzazione protegge l'inizializzazione, non ogni entità o vista.
Semplifica l'avvio, mentre gli accessi globali negli overload di comodo
rendono meno esplicite alcune dipendenze, diversamente dalla DI.

% Fonte: src/main/java/it/unifi/drivehub/AppContext.java:112-117
\noindent\begin{minipage}{\linewidth}
\begin{lstlisting}[firstnumber=112,caption={\code{AppContext.java}: accesso Singleton sincronizzato al contesto predefinito.},
  label=lst:singleton]
    public static synchronized AppContext getInstance() {
        if (instance == null) {
            instance = new AppContext();
        }
        return instance;
    }
\end{lstlisting}
\end{minipage}


\endgroup
